\chapter{含参变量积分}

\newpage

\section{含参变量常义积分}

\begin{example}
求下列极限:
\begin{enumerate}[label=(\arabic*)]
    \item $\displaystyle \lim\limits_{y\to0^+}\int_0^1\dfrac{\mathrm{d}x}{1+(1+xy)^{\frac{1}{y}}}$;
    \item $\displaystyle \lim\limits_{a\to0}\int_a^{1+a}\dfrac{\mathrm{d}x}{1+a^2+x^2}$.
\end{enumerate}
\end{example}
\exampleblank

\begin{example}
\begin{enumerate}[label=(\arabic*)]
    \item 设 $F(x)=\displaystyle\int_0^x\sin(xy)\,\mathrm{d}y$, 求 $F'(x)$;
    \item 设 $I(x)=\displaystyle\int_{\sin x}^{\cos x}e^{t^2+xt}\,\mathrm{d}t$, 求 $I'(0)$.
\end{enumerate}
\end{example}
\exampleblank

\begin{example}
\begin{enumerate}[label=(\arabic*)]
    \item 设 $f\in C(\mathbb{R}^2)$,
    \[
        F(x)=\int_0^x\left(\int_{t^2}^{x^2}f(t,s)\,\mathrm{d}s\right)\mathrm{d}t,
    \]
    求 $F'(x)$;
    \item 设 $f\in C(\lbrack0,1\rbrack^2)$,
    \[
        F(t)=\int_0^t\left(\int_0^t f(x,y)\,\mathrm{d}x\right)\mathrm{d}y,
    \]
    求 $F'(t)$.
\end{enumerate}
\end{example}
\exampleblank

\begin{example}
设
\[
    M=\int_a^b|g(x)|\,\mathrm{d}x<+\infty,
\]
$f(x,y),f_y(x,y)$ 在 $D=\lbrack a,b\rbrack\times\lbrack c,d\rbrack$ 上连续, 则
\[
    F(y)=\int_a^b g(x)f(x,y)\,\mathrm{d}x
\]
在 $\lbrack c,d\rbrack$ 上连续可微.
\end{example}
\exampleblank

\begin{example}
设 $f(x)$ 在 $\lbrack0,1\rbrack$ 上连续, 考察函数
\[
    F(t)=\int_0^1\dfrac{t}{x^2+t^2}f(x)\,\mathrm{d}x
\]
的连续性.
\end{example}
\exampleblank

\begin{example}
已知
\[
    \dfrac{1}{2\pi}\int_0^{2\pi}\dfrac{1-r^2}{1-2r\cos\theta+r^2}\,\mathrm{d}\theta=1,
    \qquad 0<r<1,
\]
试求
\[
    I(r)=\int_0^{2\pi}\ln(1-2r\cos\theta+r^2)\,\mathrm{d}\theta,
    \qquad r>0,\ r\ne1.
\]
\end{example}
\exampleblank

\begin{example}
试求
\[
    I(a)=\int_0^{\frac{\pi}{a}}\dfrac{\arctan(a\tan x)}{\tan x}\,\mathrm{d}x.
\]
\end{example}
\exampleblank

\begin{example}
试求
\[
    I(a)=\int_0^{\frac{\pi}{2}}\ln\left(\dfrac{1+a\cos x}{1-a\cos x}\right)\dfrac{\mathrm{d}x}{\cos x},
    \qquad |a|<1.
\]
\end{example}
\exampleblank

\begin{example}
试求
\[
    I(a,b)=\int_0^1\dfrac{x^b-x^a}{\ln x}\,\mathrm{d}x,
    \qquad a,b>0.
\]
\end{example}
\exampleblank

\begin{example}
设
\[
    F(r)=\int_0^{2\pi}e^{r\cos\theta}\cos(r\sin\theta)\,\mathrm{d}\theta,
\]
求证: $F(r)=2\pi$.
\end{example}
\exampleblank

\newpage

\section{含参变量反常积分}

本节基本与函数项级数讨论类似, 首先讨论一致收敛性, 再讨论相应条件下极限、求导、积分换序的应用.

\newpage

\subsection{一致收敛性}

首先需要讨论一致收敛性的判别, 这里并没有什么新的方法, 完全类似于函数项级数中的判别方法.

判别一致收敛的常用方法:
\begin{enumerate}[label=\textcircled{\arabic*},leftmargin=2.2em]
    \item 定义(包括上确界的等价定义);
    \item Cauchy 收敛准则;
    \item Weierstrass 判别法(M-判别法);
    \item Abel--Dirichlet 判别法;
    \item Dini 定理.
\end{enumerate}

判别非一致收敛的常用方法:
\begin{enumerate}[label=\textcircled{\arabic*},leftmargin=2.2em]
    \item 定义(端点法);
    \item Cauchy 收敛准则;
    \item 连续性的不可传递性.
\end{enumerate}

设
\[
    \varphi(t)=\int_a^{+\infty}f(x,t)\,\mathrm{d}x,\qquad t\in T.
\]

\begin{definition}
若对任意 $\varepsilon>0$, 存在 $A_0=A_0(\varepsilon)$, 当 $A>A_0$ 时, 有
\[
    \left|\int_A^{+\infty}f(x,t)\,\mathrm{d}x\right|<\varepsilon,
\]
则称含参变量广义积分关于 $t\in T$ 一致收敛.
\end{definition}

\begin{theorem}[\markstar\ Cauchy 收敛准则]
$\displaystyle\int_a^{+\infty}f(x,t)\,\mathrm{d}x$ 关于 $t\in T$ 一致收敛, 当且仅当对任意 $\varepsilon>0$, 存在 $A_0>a$, 当 $A,A'>A_0$ 时, 有
\[
    \left|\int_A^{A'}f(x,t)\,\mathrm{d}x\right|<\varepsilon.
\]
\end{theorem}

\begin{theorem}[\markstar\ Weierstrass 判别法]
设 $\displaystyle\int_a^{+\infty}f(x,t)\,\mathrm{d}x$ 在 $t\in T$ 上收敛. 如果
\begin{enumerate}[label=(\arabic*)]
    \item $|f(x,t)|\leq F(x)$;
    \item $\displaystyle\int_a^{+\infty}F(x)\,\mathrm{d}x$ 收敛,
\end{enumerate}
则 $\displaystyle\int_a^{+\infty}f(x,t)\,\mathrm{d}x$ 关于 $t\in T$ 一致收敛.
\end{theorem}

\begin{theorem}[\markstar\ Abel 判别法]
设
\begin{enumerate}[label=(\arabic*)]
    \item $\displaystyle\int_a^{+\infty}f(x,t)\,\mathrm{d}x$ 关于 $t\in T$ 一致收敛;
    \item $g(x,t)$ 关于 $x$ 单调, 且作为二元函数是有界的,
\end{enumerate}
则 $\displaystyle\int_a^{+\infty}f(x,t)g(x,t)\,\mathrm{d}x$ 关于 $t\in T$ 一致收敛.
\end{theorem}

\begin{theorem}[\markstar\ Dirichlet 判别法]
设
\begin{enumerate}[label=(\arabic*)]
    \item $\displaystyle\int_a^A f(x,t)\,\mathrm{d}x$ 对一切 $A\geq a,t\in T$ 一致有界;
    \item $g(x,t)$ 关于 $x$ 单调, 且 $\displaystyle\lim\limits_{x\to+\infty}g(x,t)=0$ 关于 $t\in T$ 一致,
\end{enumerate}
则 $\displaystyle\int_a^{+\infty}f(x,t)g(x,t)\,\mathrm{d}x$ 关于 $t\in T$ 一致收敛.
\end{theorem}

\begin{theorem}[\markstar\ Dini 定理]
设 $f(x,t)$ 在
\[
    D=\{(x,t)\mid a\leq x<+\infty,\ \alpha\leq t\leq\beta\}
\]
上连续且不变号,
\[
    \varphi(t)=\int_a^{+\infty}f(x,t)\,\mathrm{d}x
\]
在 $\lbrack\alpha,\beta\rbrack$ 上连续, 则 $\displaystyle\int_a^{+\infty}f(x,t)\,\mathrm{d}x$ 关于 $t\in\lbrack\alpha,\beta\rbrack$ 一致收敛.
\end{theorem}

\begin{proposition}[\markstar\ 端点法]
若存在 $y_0\in I$(或 $y_0$ 为 $I$ 的端点), 使得对任意 $A>a$, 有
\[
    \lim\limits_{y\to y_0}\int_A^{+\infty}f(x,y)\,\mathrm{d}x=B\ne0,
\]
则 $\displaystyle\int_a^{+\infty}f(x,y)\,\mathrm{d}x$ 在 $y\in I$ 上非一致收敛.
\end{proposition}

\begin{theorem}[\markstar]
若 $f(x,t)$ 在
\[
    D=\{(x,t)\mid a\leq x<+\infty,\ \alpha\leq t\leq\beta\}
\]
上连续,
\[
    \varphi(t)=\int_a^{+\infty}f(x,t)\,\mathrm{d}x
\]
在 $\lbrack\alpha,\beta\rbrack$ 上存在但不连续, 则 $\displaystyle\int_a^{+\infty}f(x,t)\,\mathrm{d}x$ 在 $\lbrack\alpha,\beta\rbrack$ 上非一致收敛.
\end{theorem}

\begin{example}
判别下列含参变量积分在特定区间的非一致收敛性:
\begin{enumerate}[label=(\arabic*)]
    \item $\displaystyle I(\alpha)=\int_0^{+\infty}e^{-\alpha x^2}\,\mathrm{d}x$, $\alpha\in(0,1)$;
    \item $\displaystyle I(\alpha)=\int_0^{+\infty}\sqrt{\alpha}\,e^{-\alpha x^2}\,\mathrm{d}x$, $\alpha\in(0,+\infty)$;
    \item $\displaystyle I(y)=\int_0^{+\infty}\dfrac{x\sin(xy)}{y(1+x^2)}\,\mathrm{d}x$, $y\in(0,+\infty)$.
\end{enumerate}
\end{example}
\exampleblank

\begin{example}
试用定义证明:
\[
    I(y)=\int_1^{+\infty}e^{-\frac{1}{y^2}(x-\frac{1}{y})^2}\,\mathrm{d}x
\]
关于 $y\in(0,1)$ 一致收敛.
\end{example}
\exampleblank

\begin{example}
证明下列命题:
\begin{enumerate}[label=(\arabic*)]
    \item $\displaystyle I=\int_1^{+\infty}\dfrac{(\ln x)^p}{x\sqrt{x}}\,\mathrm{d}x$ 关于 $p\in\lbrack0,10\rbrack$ 一致收敛;
    \item $\displaystyle I(y)=\int_1^{+\infty}\dfrac{y^2-x^2}{(x^2+y^2)^2}\,\mathrm{d}x$ 关于 $y\in\lbrack0,+\infty)$ 一致收敛;
    \item $\displaystyle \int_0^{+\infty}x\sin(x^4)\cos(\alpha x)\,\mathrm{d}x$ 对 $\alpha\in\lbrack a,b\rbrack$ 一致收敛.
\end{enumerate}
\end{example}
\exampleblank

\begin{example}
证明:
\[
    I(\alpha)=\int_0^{+\infty}\dfrac{e^{-x}}{|\sin x|^\alpha}\,\mathrm{d}x
\]
对 $\alpha\in\lbrack0,b\rbrack$, $0<b<1$, 一致收敛.
\end{example}
\exampleblank

\begin{example}
证明下列含参变量积分在对应区间上一致收敛:
\begin{enumerate}[label=(\arabic*)]
    \item $\displaystyle I(\alpha)=\int_0^{+\infty}\sin(x^\alpha)\,\mathrm{d}x$, $\alpha\geq\alpha_0$, $\alpha_0>1$;
    \item $\displaystyle I(\alpha)=\int_1^{+\infty}\dfrac{x\sin(\alpha x)}{1+x^2}\,\mathrm{d}x$, $\alpha\geq\alpha_0$, $\alpha_0>0$;
    \item $\displaystyle I(\alpha)=\int_0^{+\infty}e^{-x}\dfrac{\sin(\alpha x)}{x}\,\mathrm{d}x$, $\alpha\geq\alpha_0$, $\alpha_0>0$.
\end{enumerate}
\end{example}
\exampleblank

\begin{example}
设 $f(x)$ 在 $x>0$ 时连续, $\displaystyle\int_0^{+\infty}x^\alpha f(x)\,\mathrm{d}x$ 在 $\alpha=a$ 和 $\alpha=b$ $(a<b)$ 时收敛, 则该积分对 $\alpha\in\lbrack a,b\rbrack$ 一致收敛.
\end{example}
\exampleblank

\newpage

\subsection{极限换序}

\begin{theorem}[\markstar\ 连续性]
设 $f(x,t)$ 在
\[
    D=\{(x,t)\mid a\leq x<+\infty,\ \alpha\leq t\leq\beta\}
\]
上连续, 且 $\displaystyle\int_a^{+\infty}f(x,t)\,\mathrm{d}x$ 关于 $t$ 在 $\lbrack\alpha,\beta\rbrack$ 上一致收敛于 $\varphi(t)$, 则 $\varphi(t)$ 在 $\lbrack\alpha,\beta\rbrack$ 上连续, 即对任意 $t_0\in\lbrack\alpha,\beta\rbrack$, 有
\[
    \lim\limits_{t\to t_0}\int_a^{+\infty}f(x,t)\,\mathrm{d}x
    =\int_a^{+\infty}\lim\limits_{t\to t_0}f(x,t)\,\mathrm{d}x.
\]
\end{theorem}

\begin{theorem}[\markstar\ 可微性]
设 $f(x,t),f_t(x,t)$ 在 $D$ 上连续, $\displaystyle\int_a^{+\infty}f(x,t)\,\mathrm{d}x$ 在 $\lbrack\alpha,\beta\rbrack$ 上收敛于 $\varphi(t)$, $\displaystyle\int_a^{+\infty}f_t(x,t)\,\mathrm{d}x$ 关于 $t\in\lbrack\alpha,\beta\rbrack$ 一致收敛, 则 $\varphi(t)$ 在 $\lbrack\alpha,\beta\rbrack$ 上可微, 且
\[
    \varphi'(t)=\dfrac{\mathrm{d}}{\mathrm{d}t}\int_a^{+\infty}f(x,t)\,\mathrm{d}x
    =\int_a^{+\infty}f_t(x,t)\,\mathrm{d}x.
\]
\end{theorem}

\begin{theorem}[\markstar\ 积分换序]
条件同定理 11.2.7, 则 $\varphi(t)$ 在 $\lbrack\alpha,\beta\rbrack$ 上可积, 且
\[
    \int_\alpha^\beta\mathrm{d}t\int_a^{+\infty}f(x,t)\,\mathrm{d}x
    =\int_a^{+\infty}\mathrm{d}x\int_\alpha^\beta f(x,t)\,\mathrm{d}t.
\]
\end{theorem}

\begin{remark}
更一般地, 两个积分都是反常积分时, 换序的充分条件是 $\displaystyle\int_a^{+\infty}f(x,t)\,\mathrm{d}x$、$\displaystyle\int_\alpha^{+\infty}f(x,t)\,\mathrm{d}t$ 分别关于 $t,x$ 一致收敛, 且
\[
    \int_\alpha^{+\infty}\mathrm{d}t\int_a^{+\infty}f(x,t)\,\mathrm{d}x,
    \qquad
    \int_a^{+\infty}\mathrm{d}x\int_\alpha^{+\infty}f(x,t)\,\mathrm{d}t
\]
之一收敛即可, 具体可见基础教材 16.2 的定理 16.9 及结论.
\end{remark}

\begin{example}
讨论下列函数的连续性:
\begin{enumerate}[label=(\arabic*)]
    \item $\displaystyle F(\alpha)=\int_0^{+\infty}\dfrac{x}{2+x^\alpha}\,\mathrm{d}x$, $\alpha\in(2,+\infty)$;
    \item $\displaystyle f(\alpha)=\int_0^{+\infty}\dfrac{e^{-x}}{|\sin x|^\alpha}\,\mathrm{d}x$, $\alpha\in(0,1)$.
\end{enumerate}
\end{example}
\exampleblank

\begin{example}
证明含参变量积分
\[
    F(u)=\int_0^{+\infty}\dfrac{\sin(ux^2)}{x}\,\mathrm{d}x
\]
在 $(0,+\infty)$ 内非一致收敛, 但 $F(u)$ 在 $(0,+\infty)$ 内连续.
\end{example}
\exampleblank

\begin{example}
设 $f(x,y)$ 在 $\lbrack a,b\rbrack\times\lbrack c,+\infty)$ 上一致连续, 且
\[
    \int_c^{+\infty}|\varphi(y)|\,\mathrm{d}y
\]
收敛, 则
\[
    F(x)=\int_c^{+\infty}\varphi(y)f(x,y)\,\mathrm{d}y
\]
在 $\lbrack a,b\rbrack$ 上连续.
\end{example}
\exampleblank

\begin{example}[\markstar]
设 $f(x)$ 在 $\lbrack a,b\rbrack$ 上有界, 证明
\[
    g(\alpha)=\int_a^b\dfrac{f(x)}{\sqrt{x-\alpha}}\,\mathrm{d}x
\]
在 $\mathbb{R}$ 上连续.
\end{example}
\exampleblank

\begin{example}
证明下列命题:
\begin{enumerate}[label=(\arabic*)]
    \item $\displaystyle \lim\limits_{y\to+\infty}\int_0^y e^{-x}\,\mathrm{d}x=1$;
    \item 设 $f(x)\in C(\mathbb{R})$ 且 $f(x)\geq0$, $\displaystyle\int_0^{+\infty}f(x)\,\mathrm{d}x$ 收敛, 则
    \[
        \lim\limits_{n\to\infty}\int_0^{+\infty}n\ln\left(1+\dfrac{f(x)}{n}\right)\mathrm{d}x
        =\int_0^{+\infty}f(x)\,\mathrm{d}x.
    \]
\end{enumerate}
\end{example}
\exampleblank

\begin{example}
设 $f(x)\in R(0,+\infty)$, 则
\[
    \lim\limits_{a\to0^+}\int_0^{+\infty}e^{-ax}f(x)\,\mathrm{d}x
    =\int_0^{+\infty}f(x)\,\mathrm{d}x.
\]
\end{example}
\exampleblank

\begin{example}[\markstar]
设 $f(x)$ 是 $\lbrack0,A\rbrack$ 上单调函数, 则
\[
    \lim\limits_{\lambda\to+\infty}\int_0^A\dfrac{\sin\lambda x}{x}f(x)\,\mathrm{d}x
    =\dfrac{\pi}{2}f(0^+).
\]
\end{example}
\exampleblank

\begin{example}
计算积分
\[
    g(\alpha)=\int_1^{+\infty}\dfrac{\arctan(\alpha x)}{x^2\sqrt{x^2-1}}\,\mathrm{d}x.
\]
\end{example}
\exampleblank

\begin{example}
计算积分
\[
    g(\alpha,\beta)=\int_0^{+\infty}\dfrac{\arctan(\alpha x)\arctan(\beta x)}{x^2}\,\mathrm{d}x.
\]
\end{example}
\exampleblank

\begin{example}
计算积分
\[
    I=\int_0^{+\infty}\dfrac{e^{-\alpha x}-e^{-\beta x}}{x}\cos(mx)\,\mathrm{d}x,
    \qquad \alpha,\beta>0.
\]
\end{example}
\exampleblank

\begin{example}
设
\[
    f(x)=\left(\int_0^x e^{-t^2}\,\mathrm{d}t\right)^2,
    \qquad
    g(x)=\int_0^1\dfrac{e^{-x^2(1+t^2)}}{1+t^2}\,\mathrm{d}t.
\]
证明:
\begin{enumerate}[label=(\arabic*)]
    \item $f(x)+g(x)\equiv C$, 并确定此常数;
    \item $\displaystyle\int_0^{+\infty}e^{-t^2}\,\mathrm{d}t=\dfrac{\sqrt{\pi}}{2}$.
\end{enumerate}
\end{example}
\exampleblank

\begin{example}
利用
\[
    \dfrac{1}{\sqrt{t}}=\dfrac{1}{\sqrt{\pi}}\int_0^{+\infty}e^{-tu^2}\,\mathrm{d}u
\]
计算 Fresnel 积分
\[
    \int_0^{+\infty}\sin x^2\,\mathrm{d}x.
\]
\end{example}
\exampleblank

\begin{example}
试利用
\[
    \int_0^{+\infty}e^{-u^2}\,\mathrm{d}u
    \overset{u=\alpha x}{=}
    \int_0^{+\infty}\alpha e^{-\alpha^2x^2}\,\mathrm{d}x,
    \qquad \alpha>0,
\]
计算
\[
    \int_0^{+\infty}e^{-x^2}\,\mathrm{d}x.
\]
\end{example}
\exampleblank

\begin{example}
若
\[
    f(x)=\sum\limits_{n=0}^{\infty}a_nx^n,\qquad a_n>0,
\]
的收敛半径为 $+\infty$, 且 $\displaystyle\sum\limits_{n=0}^{\infty}a_nn!$ 收敛, 则
\[
    \int_0^{+\infty}e^{-x}f(x)\,\mathrm{d}x
\]
收敛, 且
\[
    \int_0^{+\infty}e^{-x}f(x)\,\mathrm{d}x
    =\sum\limits_{n=0}^{\infty}a_nn!.
\]
\end{example}
\exampleblank

\begin{example}
试用级数法求
\[
    f(x)=\int_0^{\frac{\pi}{2}}\ln(1-x^2\cos^2\theta)\,\mathrm{d}\theta,
    \qquad |x|<1.
\]
\end{example}
\exampleblank

\begin{example}
计算积分
\[
    \int_0^{+\infty}\dfrac{x}{1+e^x}\,\mathrm{d}x.
\]
\end{example}
\exampleblank

\begin{example}[\markstar]
试证
\[
    \int_0^{+\infty}\dfrac{x^{s-1}}{e^x+1}\,\mathrm{d}x
    =\Gamma(s)(1-2^{1-s})\xi(s),\qquad s>1,
\]
其中
\[
    \xi(s)=\sum\limits_{n=1}^{\infty}\dfrac{1}{n^s}.
\]
\end{example}
\exampleblank

\newpage

\section{\texorpdfstring{Euler 积分——$B$ 函数与 $\Gamma$ 函数}{Euler 积分——B 函数与 Gamma 函数}}

\begin{proposition}[\markstar]
$\Gamma$ 函数与 $B$ 函数基本性质如下.

\begin{enumerate}[label=(\arabic*)]
    \item 等价形式:
    \begin{align*}
        \Gamma(s)
        &=\int_0^{+\infty}x^{s-1}e^{-x}\,\mathrm{d}x
        =2\int_0^{+\infty}t^{2s-1}e^{-t^2}\,\mathrm{d}t\\
        &=\int_0^1\left(\ln\dfrac{1}{t}\right)^{s-1}\,\mathrm{d}t
        =\lim\limits_{n\to\infty}\dfrac{n!\,n^s}{s(s+1)\cdots(s+n)},\qquad s>0,
    \end{align*}
    且
    \begin{align*}
        B(p,q)
        &=\int_0^1x^{p-1}(1-x)^{q-1}\,\mathrm{d}x\\
        &=2\int_0^{\frac{\pi}{2}}\cos^{2p-1}\theta\sin^{2q-1}\theta\,\mathrm{d}\theta\\
        &=\int_0^{+\infty}\dfrac{u^{p-1}}{(1+u)^{p+q}}\,\mathrm{d}u\\
        &=\int_0^1\dfrac{u^{p-1}+u^{q-1}}{(1+u)^{p+q}}\,\mathrm{d}u,
        \qquad p,q>0.
    \end{align*}

    \item 递推公式:
    \[
        B(p+1,q)=\dfrac{p}{p+q}B(p,q),\qquad
        B(p,q+1)=\dfrac{q}{p+q}B(p,q),\qquad
        \Gamma(p+1)=p\Gamma(p).
    \]
    并且
    \[
        B(1,1)=\Gamma(1)=1,\qquad
        \Gamma\left(\dfrac{1}{2}\right)=\sqrt{\pi},\qquad
        \Gamma(n)=(n-1)!,\quad n\in\mathbb{Z}^+.
    \]
    对正整数 $m,n$, 有
    \[
        B(n,m)=\dfrac{(n-1)!(m-1)!}{(m+n-1)!}.
    \]

    \item $\Gamma$ 函数与 $B$ 函数关系:
    \[
        B(p,q)=\dfrac{\Gamma(p)\Gamma(q)}{\Gamma(p+q)}.
    \]

    \item 余元公式:
    \[
        B(p,1-p)=\Gamma(p)\Gamma(1-p)=\dfrac{\pi}{\sin p\pi}.
    \]

    \item 倍元公式:
    \[
        \Gamma(2x)=\dfrac{2^{2x-1}}{\sqrt{\pi}}\Gamma(x)\Gamma\left(x+\dfrac{1}{2}\right).
    \]

    \item 求导公式: $\Gamma(x)$ 在 $(0,+\infty)$ 内闭一致收敛, 且有连续的各阶导数,
    \[
        \Gamma^{(n)}(x)=\int_0^{+\infty}t^{x-1}e^{-t}(\ln t)^n\,\mathrm{d}t.
    \]
\end{enumerate}
\end{proposition}

\begin{example}
计算下列积分:
\begin{enumerate}[label=(\arabic*)]
    \item $\displaystyle\int_a^b(x-a)^p(b-x)^q\,\mathrm{d}x$, $p,q>-1$, $b>a$;
    \item $\displaystyle\int_0^{+\infty}\dfrac{\mathrm{d}x}{1+x^3}$;
    \item $\displaystyle\int_0^{+\infty}\dfrac{\mathrm{d}x}{1+x^4}$.
\end{enumerate}
\end{example}
\exampleblank

\begin{example}
证明下列命题:
\begin{enumerate}[label=(\arabic*)]
    \item
    \[
        \int_0^{+\infty}\dfrac{x^{p-1}}{(a+bx^q)^r}\,\mathrm{d}x
        =\dfrac{a^{\frac{p}{q}-r}}{q\,b^{\frac{p}{q}}}
        B\left(\dfrac{p}{q},r-\dfrac{p}{q}\right),\qquad a,b>0;
    \]
    \item
    \[
        \int_0^{+\infty}\dfrac{\mathrm{d}x}{(x^2+1)^n}
        =\dfrac{\sqrt{\pi}\,\Gamma\left(n-\dfrac{1}{2}\right)}{2\Gamma(n)}
        =\dfrac{(2n-3)!!}{(2n-2)!!}\dfrac{\pi}{2}.
    \]
\end{enumerate}
\end{example}
\exampleblank

\begin{example}
计算下列积分:
\begin{enumerate}[label=(\arabic*)]
    \item $\displaystyle\int_0^{\frac{\pi}{2}}\sin^\alpha x\,\mathrm{d}x$, $\alpha>0$;
    \item $\displaystyle\int_0^{\frac{\pi}{2}}\sin^\alpha x\sin^\beta x\,\mathrm{d}x$;
    \item $\displaystyle\int_0^{\frac{\pi}{2}}\tan^\alpha x\,\mathrm{d}x$, $|\alpha|<1$;
    \item $\displaystyle\int_0^\pi\dfrac{\mathrm{d}x}{\sqrt{3-\cos x}}$;
    \item $\displaystyle\int_0^\pi\dfrac{\sin^{n-1}x}{(1+k\cos x)^n}\,\mathrm{d}x$, $0<|k|<1$.
\end{enumerate}
\end{example}
\exampleblank

\begin{example}
计算下列积分:
\begin{enumerate}[label=(\arabic*)]
    \item $\displaystyle\int_0^{+\infty}\dfrac{x^{p-1}\ln x}{1+x}\,\mathrm{d}x$;
    \item $\displaystyle\int_0^{+\infty}\dfrac{x\ln x}{1+x^2}\,\mathrm{d}x$;
    \item $\displaystyle\int_0^{+\infty}\dfrac{\ln^2x}{1+x^4}\,\mathrm{d}x$.
\end{enumerate}
\end{example}
\exampleblank

\begin{example}
计算积分:
\begin{enumerate}[label=(\arabic*)]
    \item $\displaystyle\int_0^1\ln\Gamma(x)\,\mathrm{d}x$;
    \item $\displaystyle\int_0^1\ln\Gamma(x)\sin\pi x\,\mathrm{d}x$.
\end{enumerate}
\end{example}
\exampleblank

\begin{example}[\markstar]
应用 $\Gamma$ 函数的 Gauss 乘积分解公式
\[
    \Gamma(x)=\lim\limits_{n\to\infty}\dfrac{n^x\cdot n!}{x(x+1)\cdots(x+n)},
\]
证明:
\begin{enumerate}[label=(\arabic*)]
    \item
    \[
        \dfrac{\Gamma'(x)}{\Gamma(x)}
        =-\gamma+\sum\limits_{n=1}^{\infty}\left(\dfrac{1}{n}-\dfrac{1}{x+n-1}\right),
    \]
    其中 $\gamma$ 为 Euler 常数;
    \item $-\Gamma'(1)=\gamma$.
\end{enumerate}
\end{example}
\exampleblank

\begin{example}[\markstar]
设 $t>1$, 证明
\[
    \int_1^{+\infty}\dfrac{(\ln x)^{t-1}}{x(x-1)}\,\mathrm{d}x
    =\Gamma(t)\xi(t),
\]
其中
\[
    \xi(t)=\sum\limits_{n=1}^{\infty}\dfrac{1}{n^t}
\]
为 Zeta 函数.
\end{example}
\exampleblank
